\documentclass[../main]{subfiles}
\begin{document}
\ifSubfilesClassLoaded{\setcounter{page}{19}}{}
\section{Комунізм як реальність у суспільному виробництві}
\label{sec:04_kommunizm_realnost}

Твердження про те, що комунізм вже існує, може здатися надто сміливим. У будь-якому разі, воно потребує підтвердження. Навпаки, одразу впадає в око відсутність комунізму, панування праці, цінність якого для того, хто працює, визначається лише у формі оплати праці – заробітної плати. Може здатися, що зародкова форма комунізму – безпосередньо колективна форма вільної людської діяльності – є чимось, що не стоїть на головній дорозі розвитку людського суспільства. Тим часом, ми спробуємо показати, що це не так. Для цього недостатньо підібрати лише ті приклади, які б засвідчували наявність комунізму в обмежених рамках експлуататорського класового суспільства. Необхідно розглянути безпосередньо колективну форму вільної діяльності, що здійснюється в сучасному класовому суспільстві, з точки зору можливостей її подальшого розвитку.

Нагадаємо, свого часу Ленін зміг по-новому поглянути на радянську владу, а згодом і на комуністичні суботники. «Особливе, навіть гігантське значення в цьому плані має те, що робітники організовують комуністичні суботники з власної ініціативи. Мабуть, це лише початок, але початок надзвичайно важливий. Це – початок перевороту, який є більш складним, більш важливим, більш докорінним і вирішальним, ніж повалення буржуазії, оскільки це – перемога над власною косністю, розхлябаністю, дрібнобуржуазним егоїзмом, над тими звичками, які проклятий капіталізм залишив у спадок робітникам і селянам». Коли ця перемога буде закріплена, тоді й лише тоді буде створено нову суспільну дисципліну, соціалістичну дисципліну, тоді й лише тоді повернення назад, до капіталізму, стане неможливим, і комунізм стане справді непереможним». Ленін стверджував, що перший комуністичний суботник, організований 10 травня 1919 року залізничними робітниками Московсько-Казанської залізниці, має більше історичне значення, ніж будь-яка перемога в Першій світовій війні. Йшлося не просто про самостійне волевиявлення робітників, які організували суботник, а й про суттєве збільшення продуктивності праці в цій ініціативі, зумовлене тим, що для кожного працівника це була його особиста справа, де реалізувався особистий інтерес у розвитку залізничного господарства в інтересах усього суспільства.

Після публікації розділу про заснування комунізму один із наших читачів зробив важливе зауваження: «Комуністичний спосіб виробництва передбачає не лише виробництво матеріальних благ, а не просто пушкінський літературний гурток, який сприяє становленню поетів. Адже сьогодні повно так званих «комуністичних відносин», заснованих на експлуатації для «обраних», – стверджував читач. Він вимагав надати докази того, що комунізм вже існує і яким він може бути саме у сфері матеріального виробництва. На нашу думку, правота читача полягає у двох аспектах. По-перше, читач виступив від імені пролетаріату – людей, які безпосередньо зайняті в матеріальному виробництві шляхом продажу власної робочої сили, де комунізму не спостерігається. По-друге, читач порушує надзвичайно важливе питання щодо комуністичного матеріального виробництва, адже саме воно визначає комунізм у суспільному масштабі, а не прагнення, побажання чи навіть дії політичних революціонерів.

Проте, комунізм неможливий у чисто матеріальному виробництві. Там, де людина безпосередньо виробляє лише речі, а виробництво суспільних відносин відбувається в цьому процесі цілком незалежно від неї, від її цілей, волі та свідомості, не може йтися про жодний комунізм, навіть якщо саме по собі виробництво речей приносить працівнику певне задоволення від праці. У такому виробництві людина є лише матеріальною силою поряд з іншими матеріальними силами. Зрештою, вона виступає як річ, що впливає на інші речі. З погляду марксизму, такий вид праці для настання комунізму не лише може, але й повинен бути знищений як людська праця і повністю переданий машинам. Говорячи про історичну роль капіталу, Маркс особливо наголошував на важливості створення додаткової вартості та універсальності людських потреб як умови і передумови комунізму. З цим капітал вже справляється і за часів капіталізму. Але лише цього недостатньо. Історична роль капіталу буде остаточно виконана, а комуністичне суспільство стане реальністю лише тоді, «коли припиниться така праця, коли людина сама робитиме те, що вона може змусити речі робити для себе, для людини». У цьому сенсі людина не може діяти по-комуністичному, поки її не виведуть за рамки чисто матеріального виробництва, поки вона робить те, що можуть робити речі, і, що важливо, у цьому процесі сама виступає як річ. Тому пролетаріат – єдиний клас, який до кінця послідовно зацікавлений у тому, щоб історична роль капіталу була виконана до кінця. Це і є знищення тих умов, за яких суспільство постійно розділяється на класи.

Автор зауваження слушно зазначив, що у сфері виробництва ідеального комунізму можна знайти безліч прикладів. Сюди можна віднести роботу вченого та колективу вчених, які, щоб дати людям наукову істину, працюють і вдень, і вночі, а ще й читають просвітницькі лекції у вільний час. Певний мінімум матеріальних благ тут є умовою, а не метою комуністичної діяльності. Комуністичною діяльністю є робота письменника, художника, музиканта, які працюють для того, щоб люди бачили, чули, відчували і розуміли світ краще, щоб розширити межі людських можливостей, а не заради грошей чи слави тощо. Звісно, далеко не вся діяльність у сфері ідеального є комуністичною. А якщо згадані в попередній главі моменти безпосередньо-колективної форми вільної людської діяльності очевидні, то ми маємо справу з локальним комунізмом, хоча тут комунізм існує у досить обмеженій формі. Тут, як справедливо зазначено, комунізм далеко не в повній мірі реалізований.

По-перше, самі умови для його існування створюються за рахунок чисто капіталістичної експлуатації людини людиною. По-друге, капітал, як домінуюча суспільна формація, і тут обмежує свободу жорсткими рамками. Не можна побудувати комунізм як спосіб матеріального виробництва, не виходячи за межі сфери ідеального. І хоча комунізм у сфері ідеального протягом усієї історії людства відігравав найважливішу роль, оскільки у вільній діяльності народжувалися ті ідеї, які не просто відображають, а й «творять» світ, все ж він не може відігравати вирішального значення, оскільки ще не робить суспільство комуністичним. По-третє, тут поки що не йдеться про подолання суспільного поділу праці, що є найважливішою умовою для комунізму. У цьому контексті всебічність людини може проявлятися лише на рівні свідомості, але не на рівні цілісності людини як суб’єкта матеріально-практичної діяльності.

Проте існує таке матеріальне виробництво, яке не є чисто матеріальним виробництвом, а є виробництвом людини в її цілісності. Таке виробництво не просто опосередковано визначає, а включає в себе виробництво ідеального – суспільної свідомості, суспільного ідеалу, суспільного цілепокладання. Наведений вище приклад із комуністичним суботником, під час якого робітники займалися навантаженням і розвантаженням вантажів, а також ремонтом рухомого складу, добре це показує. Неприваблива, на перший погляд, ручна праця тут була не просто працею з перенесення вантажів, а працею з виробництва нової соціальної реальності. Тут доречно згадати про суботники з прибирання цукрової тростини на Кубі: За допомогою досить примітивних знарядь праці, під палючим сонцем, учасники не просто рубали очерет, а здобували кошти для модернізації виробництва, будівництва шкіл, лікарень – заради створення умов для розвитку всіх представників кубинського народу. Така праця була і в колонії імені Горького та комуні імені Дзержинського, якими керував Антон Семенович Макаренко. Будь-яка справа, навіть викопування за допомогою лопати вигрібних ям і зведення над ними дерев’яних будиночків, ставала тут, зрештою, процесом створення нової людини, вільною людською діяльністю, що відповідає критеріям істини, добра і краси.

Отже, питання про вільну форму безпосередньо-колективної діяльності в матеріальному виробництві слід поставити так: який зв’язок із суспільством, як цілісним організмом у його розвитку, ця діяльність втілює та створює для тих, хто працює, і яке місце та роль вона відіграє в розвитку суспільного цілого?

Легко помітити, що наведені вище приклади взяті з історії соціалізму. При соціалізмі, самим фактом формального усуспільнення засобів виробництва, усуваються певні перешкоди для вільного розвитку суспільного життя. Тому хтось може зауважити, що ці приклади не зовсім коректні. Але, по-перше, усунення бар’єрів дає можливість реалізації не того, чого взагалі не існує, і не того, для чого взагалі немає умов, а того, для чого не було повного комплексу умов, і єдине, чого не вистачало в цьому комплексі – це усунення панування капіталу. Те, що усунення бар’єрів дає простір для розвитку комуністичних тенденцій, означає, як мінімум, що ці тенденції існують, що потенціал для становлення такої діяльності вже криється в капіталізмі.

Нам також можуть заперечити, що суботник – це щось нестійке, ситуативне, що не визначає загальний хід життя. Хоча в процесі суботника ставляться і досягаються як суспільні цілі, так і особисті цілі учасників суботника, за межами суботника в матеріальному виробництві діють ті ж самі капіталістичні закони: робоча сила є товаром, і кожен окремий пролетарій укладає договір купівлі-продажу робочої сили зі своїм же класом, який виступає по відношенню до нього як колективний капіталіст; закон вартості зберігає свою силу, якщо його дію спеціально не обмежувати, доводячи тенденцію капіталу до автоматизації виробництва до межі.

Отже, наведений вище приклад сам по собі ще не перетворює суспільство на комуністичне. Ці зауваження цілком обґрунтовані, і ми з ними не сперечатимемося. Але суть полягає не лише і не стільки в простому фіксуванні, мовляв, це – комунізм, а це – ні, або не зовсім комунізм, а в тих тенденціях суспільства, потенціал яких проявляється у таких «острівцях комунізму». Саме в цьому дусі міркував Ленін про «паростки комунізму». Суботники цікавили його саме як паростки, зародки. Тому він звертає увагу, передусім, на дві речі: 1) вільне вираження волі учасників суботника щодо засобів виробництва і щодо самих себе; 2) значне зростання продуктивності праці під час суботників, яке досягається завдяки цьому вільному вираженню волі.

Ще одним важливим зауваженням було б те, що в наведених вище прикладах діяльність здійснюється далеко не найсучаснішими інструментами, навіть для капіталізму, і, тим більше, не на комуністичній технологічній базі. Це напруження людських сил відбувається радше через недостатній рівень суспільного розвитку: брак техніки, недостатня організація технологічних процесів, навіть у чисто капіталістичному плані. На нашу думку, саме в цьому полягає найобґрунтованіше заперечення проти такого комунізму, як зародка нового суспільства. Комунізм, що виник із потреби, про який йшлося вище, ще не є комунізмом за внутрішньою необхідністю, і тому з нього самого ще не випливає розвиток комуністичного суспільства. Він є, можливо, і суттєвою, але лише передумовою для такого суспільства. Тут, знову ж таки, необхідно дослідити зв’язок і наступність у способах побудови колективу між таким типом комунізму та комунізмом, який розвивається на власній основі, і цьому потрібно присвятити окрему працю.

Однак, посилання на необхідність окремої розробки цієї проблеми зовсім не звільняє нас від обов’язку вказати на практичну можливість цього зв’язку в загальних рисах, а також від обов’язку надати свідчення існування безпосередньо-колективних форм вільної людської діяльності, заснованих на передових сучасних технологіях.

Зв’язок між «комунізмом, зумовленим потребою», та комунізмом, що має власну основу, їхня наступність може полягати, по-перше, у вихованні мас, у формуванні звички до вільного волевиявлення відповідно до суспільної потреби у виробництві суспільного життя; по-друге, у створенні комуністичних форм колективної організації, які в умовах подолання суспільної потреби, що їх і породила, зовсім не обов’язково повинні зникнути, а з урахуванням виховання, яке відбувається в них у дусі справжньої свободи, як усвідомленої суспільної необхідності, можуть стати формами комуністичної організації суспільного життя. Ми не утопісти. Як саме це буде здійснюватися, ми не будемо вигадувати. Це щоразу залежить від переплетення конкретно-історичних обставин. Ми лише констатуємо можливість такого подальшого розвитку, що полягає в подібних проявах колективно-організаційної форми людської діяльності, яка передбачає вільну та рівну участь у ній.

Що стосується комунізму, заснованого на передових технологіях, то подібні прояви ми бачимо в процесі спільного використання сучасних інформаційних технологій та їх спільного створення. Наприклад, операційна система Linux та інші програми, створені в рамках розробки вільного програмного забезпечення, створені на комуністичних засадах.

Уточнімо, що мова йде не про всі програми з відкритим кодом. Багато великих сучасних проєктів стали основою для величезних корпорацій, і те, що код є відкритим, не означає, що діяльність зі створення програм є вільною. І там, де це має місце, корпорація повністю визначає політику проєкту, фінансує його та отримує від цього прибуток. Код пишуть програмісти за зарплату. Тут немає нічого спільного з комунізмом (хоча існує можливість перетворити таку діяльність на комуністичну). Наприклад, так створюються допоміжні продукти, які потрібні одразу кільком компаніям. Кооперація здешевлює розробку. Але зрештою, продукт тут створюється для отримання прибутку. Часто для будь-якого використання такого продукту, а тим більше для комерційного, потрібні великі ресурси, які є лише у великих корпорацій. У них зазвичай є комерційний сервіс або програмне забезпечення, для якого розроблене є однією з умовних «частин», тобто умовою для його роботи та успішних продажів. Але корпорації монополізували те, що технологічно має значний потенціал, який виходить за межі капіталізму. Цей потенціал пов’язаний з технічною можливістю як виробництва, так і використання програмного забезпечення завдяки розвитку інформаційних машин і мереж.

Усе почалося з вільного програмного забезпечення, і певна частина проєктів зараз існує саме як проєкти, де немає товарно-грошових відносин. Більше того, розробники ведуть боротьбу за те, щоб ніхто не отримував з цього комерційної вигоди. Це і є емпірично даний комунізм у рамках капіталізму. Відкрите програмне забезпечення передбачає повний вільний доступ до програм і до коду програм, а також надання у вільний доступ самих програм та їхнього коду, створених з використанням отриманих раніше вільних вихідних кодів. Наприклад, операційна система Linux виявилася конкурентоспроможною щодо комерційної Windows. Корпорації зрештою намагаються підкорити її і створюють продукти на основі відкритого коду цієї системи. Але сама вона, як і раніше, залишається безкоштовною.

Правда, те, що створюється в рамках некомерційного вільного програмного забезпечення, далеко не завжди працює так, як було задумано, і користувачі на це звертають увагу. Але й комерційне програмне забезпечення на шляху до стадії повністю готового продукту проходить через це, а деякі розробки взагалі закриваються через їхню безперспективність. Те саме відбувається і тут. Тільки тут усі етапи відкриті та видимі для всіх. Незавершені та безперспективні проєкти також стають надбанням громадськості. Вільне програмне забезпечення – це взагалі часто напівфабрикат, розрахований на те, щоб користувач довів його до стану повної готовності. Це ускладнює його конкурентоспроможність в умовах капіталізму у світі із загальним низьким рівнем інформаційної грамотності.

Корпорації розробляють свій продукт, враховуючи цей рівень обізнаності, доводячи його до можливості примітивно-інтуїтивного використання. Їхній продукт призначений для людей, яким в умовах відчуження легше заплатити, щоб користування було інтуїтивним, ніж витрачати час на те, щоб розібратися в програмі. Продукт корпорацій не передбачає, що користувачі створюватимуть нові програми або модифікуватимуть існуючі під себе шляхом доповнення або переписування частини коду. Його використання не вимагає особливо високої культури і тим більше розуміння технологічних процесів. Тому програми, вироблені великими корпораціями, мають набагато ширші можливості для збуту. Крім того, вони часто мають вищу якість, оскільки опрацьовані в усіх деталях.

На противагу цьому, безкоштовне програмне забезпечення, яке не передбачає комерційної вигоди, вимагає хоча б базових знань програмування, передбачає можливість, а часто й потребує доопрацювання та переписування коду. Воно змушує користувача частково ставати програмістом і постійно вдосконалювати свої навички в цьому \footnote{Звісно, це може дратувати і заважати зосередитися на виконанні завдання. В результаті, замість того, щоб робити справу, людям доводиться постійно лагодити несправний інструмент для Linux. Але немає жодних принципових перешкод (якщо дивитися з точки зору потенціалу) для подолання цього і написання більш якісних програм, за умови значних витрат суспільного часу на їх створення.}. Саме цей фактор, який обмежує чисто капіталістичну конкурентоспроможність безкоштовного програмного забезпечення, свідчить про його більший суспільний потенціал у порівнянні з комерційним (мова йде не про відмову від зручного інтерфейсу, а про можливість самостійно створювати зручний інтерфейс). У його розробці, способах передачі та використанні закладено не просто потенціал технічного розвитку, а й потенціал розвитку людини – подолання поділу праці, підвищення рівня всіх учасників процесу до рівня сучасної комп’ютерної грамотності. І хоча цей локальний комунізм намагаються (і іноді успішно) підкорити собі капіталістичні компанії, його цілком можна вважати паростком нового суспільства, хоч поки що слабким і тендітним. По-перше, він виник на основі найпередовіших для капіталізму технологій. По-друге, він чітко окреслює межу між капіталістичним способом використання та створення цих технологій і способом, що виходить за рамки капіталізму.

Важливо також те, що таке виробництво вже саме по собі, з точки зору технологій, веде до руйнування ієрархії у виробництві та розмежування між користувачами і програмістами. Воно не потребує зовнішнього зв’язку між ними та спеціальних суспільних «органів управління», що забезпечують цей зв’язок. Зовнішнє управління людьми, навіть у найчистішій демократичній формі – у формі вирішення питань шляхом загального голосування – тут вже не потрібне.

У тенденції воно відмирає, стаючи надлишковим. Нам можуть зауважити, що, коли ми говоримо про некомерційне вільне виробництво програмного забезпечення, знову йдеться про сферу виробництва ідеального. Дійсно, процес створення коду програм, як різновид тексту, належить до ідеальної діяльності, а код програми (але не сама програма в її звичайній роботі) є ідеальним продуктом. Однак, це та ідеальна діяльність, яка є найважливішою частиною сучасного промислового матеріального виробництва. Створення програмного забезпечення – це праця з впровадження автоматизації в промислове виробництво. І, хоча код програми є ідеальним продуктом, функціонування сучасних промислових машин на основі цих програм (тобто «робота» програм у промисловості) є чисто матеріальним машинним виробництвом.

Звісно, у рамках розробки вільного програмного забезпечення створюється безліч безглуздих і непотрібних програм. Але це не надто важливо. Важливо те, що створюються і можуть створюватися не лише вони, а й цілком придатні для використання в масштабах сучасної індустрії продукти. І ці продукти створюються на комуністичних засадах. Хоча часто ці проєкти або не завершуються, або затягуються на роки через те, що учасники повинні заробляти на життя, а отже, в них не так багато часу на вільну діяльність.

Те саме ми спостерігаємо в процесі створення та використання сучасних баз даних. Знову ж таки, далеко не всі вони створюються на комуністичних засадах. Але існують великі загальнодоступні бази даних, які передбачають виключно спільне створення та використання. Інакше їх просто неможливо було б створити. В іншому випадку вони втрачали б сенс і ставали непридатними для тих цілей, заради яких і були створені. Наприклад, деякі просторові бази даних передбачають вільну участь у їх створенні для формування уявлення про простір в інтересах розвитку суспільства.

Тут також об’єктивні закономірності людської діяльності в тенденції виключають обов’язкову ієрархічність, притаманну капіталістичному способу виробництва, включаючи до процесу всіх учасників як рівних і вільних. Структура бази даних також є ідеальним продуктом, – заперечать нам.

Але під час створення баз даних ідеальна діяльність настільки тісно пов’язана з матеріальним виробництвом (в тому числі й матеріальним відображенням за допомогою баз) і відтворенням суспільного життя, що, попри всю її ідеалізованість, її навряд чи можна віднести до надбудової. Як і у випадку з програмним забезпеченням, це ідеальний момент матеріального суспільного виробництва, який не є абстрактним (не в свідомості, а реально) і не відокремлений від цього виробництва у вигляді окремих форм суспільної свідомості.

Отже, тут ми вказуємо не просто на ідеальну діяльність, а на промислове матеріальне виробництво, і на тенденції та потенціал, що містяться в ньому вже за часів капіталізму, які, після усунення бар’єрів у вигляді приватної власності на засоби виробництва, можуть розвинутися в комунізм, створивши як чисто технологічну, так і колективно-організаційну, виховну базу, що формує нову людину.

Подальший розвиток цієї тенденції веде до автоматизованої фабрики, в якій Маркс побачив прообраз майбутнього матеріального виробництва. Тут також закладено потенціал реалізації ідей В. М. Глушкова про «загальнодержавні» (всезагальні) системи управління економікою.

(цією автоматизованою фабрикою в суспільному масштабі). У перспективі (за умови усунення соціальних бар’єрів) кожен може мати доступ до всієї інформації про стан усього суспільного виробництва в режимі реального часу. Кожен може брати участь в управлінні не просто виробництвом речей (виробництво речей має сенс лише в тій мірі, в якій воно слугує головній меті виробництва), а виробництвом людини – суспільним життям у його цілісності.

Цікаво, що участь у всіх цих процесах передбачає виробництво та споживання ідеального продукту всіма учасниками, але такого, який не відчужений і не може бути відчужений від матеріального виробництва. У цьому простежується виникнення нового типу суспільної свідомості, що усуває його відчужені форми.

І потреба в усуненні соціальних бар’єрів для вільної діяльності також створюється сучасним капіталізмом. Сам він, незалежно від чиєїсь волі та свідомості, об’єктивно ставить це питання на порядок денний. Об’єктивна логіка розвитку сучасного капіталу, який постійно породжує війни, тим самим унеможливлює звичні способи та форми діяльності людей. Капітал змушує людей створювати якісно нові форми діяльності для вирішення проблем відтворення суспільного життя в протидії руйнівній силі капіталу.
\end{document}
